\section{数据说明与处理}
\subsection{数据来源}
附件1为问题一提供单日电价、负载和光伏预测；附件2--4分别提供全年实际供需、定时光伏预报与波动电价，用于连续调度。各附件的数据规模见表\ref{tab:data}。
\begin{table}[H]
\centering
\caption{附件数据内容与规模}
\label{tab:data}
\small
\begin{tabular}{lp{7.4cm}l}
\toprule
附件 & 数据内容 & 数据规模 \\
\midrule
1 & 单日电价、负载功率、光伏预测功率 & 144个时段 \\
2 & 2025年实际负载、实际光伏，各存一张表 & 每表$365\times144$ \\
3 & 每日0、6、12、18时发布的未来1--24小时光伏预报 & $1460\times24$个值 \\
4 & 2025年十分钟波动电价 & $365\times144$ \\
\bottomrule
\end{tabular}
\end{table}
附件1未给出观测日期，故问题一按日内时刻组织输入，不将计算用参考日期视为实际观测日期。原始附件保留不变，后续解析与绘图均使用副本。

\subsection{时间对齐与电量换算}
附件1的时间列同时含Excel数值时间和字符串时间，首个标签为0:10，末个为\texttt{0:00+1}。将标签统一换算为相对日初的分钟数，再按右端点对齐：0:10对应$[0{:}00,0{:}10)$，\texttt{0:00+1}对应$[23{:}50,24{:}00)$。设日初为$\tau_0$，则
\begin{equation}
  \tau_k=\tau_0+k\Delta t,\qquad
  I_k=[\tau_k,\tau_{k+1}),\quad \Delta t=\frac16\ \mathrm h.
  \label{eq:timegrid}
\end{equation}
一天包含144个行动区间和145个储能状态边界。负载与光伏的单位为kW，换算为区间电量后有
\begin{equation}
  \ell_k=P_k^{\mathrm L}\Delta t=\frac{P_k^{\mathrm L}}6,
  \qquad v_k=P_k^{\mathrm{PV}}\Delta t=\frac{P_k^{\mathrm{PV}}}6.
  \label{eq:power-energy}
\end{equation}
电价单位为元/kWh，区间购电费为$p_kg_k$，无需再乘区间时长。

\subsection{完整性检查与预报组织}
问题一输入完整覆盖144个区间，电价、负载与光伏均为有限数值，负载和光伏非负。计算得全天负载电量111\,024.808 kWh、光伏预测电量55\,482.836 kWh，输入曲线见图\ref{fig:input}。光伏在日间出现富余，而晚间负载主要依赖购电与储能供给；这使储能既承担光伏转移，也承担电价差利用。
\begin{figure}[H]
  \centering
  \PaperGraphic[width=0.96\linewidth]{figures/q1_input.pdf}
  \caption{附件1的日内负载、光伏预测与电价}
  \label{fig:input}
\end{figure}

附件3按日期块组织预报，同一天后续行的空日期仅在该日四行内继承，不对功率、价格或预测值填充。每个预报值同时保留发布时间和目标时刻，组合或重采样时以目标时刻对齐，以发布时间判断决策时是否可用。问题四的具体预报组合与插值方法见第\ref{sec:q4}节。

\subsection{结果汇总与模板映射}
正文购电量按式\eqref{eq:timegrid}的实际日内区间提取，四小时充放电量由连续24个区间相加，初末储电量取对应边界状态。附件5计划表的首末标签与内部网格存在差异，输出时保留模板原标签，按行序写入144个区间；这一映射不改变正文所述时段，也不视为对官方标签的勘误。
